\section{Explicit hulls and separation on structured blocks}\label{sec:structured-oracles}
The extended formulations in \cref{sec:compression} apply to arbitrary
graphs. On some block classes, their state coordinates can be projected
out explicitly while preserving unit flow and product coefficients. We
first give a linear-arithmetic oracle for cycles and theta blocks, then a
transportation oracle for any number of parallel paths. All graph classes
refer to the underlying multigraph: orientations remain arbitrary, and
different blocks may meet at articulation vertices. Bridges retain
\eqref{eq:bridge-product}; independent loops use the scalar cycle formulas.

\subsection{Cycle and theta state domains}\label{subsec:theta-domains}
A theta block consists of three internally vertex-disjoint paths between
two terminals. Orient their traversals from the first terminal to the
second. The three circulation deviations sum to zero, so write them as
$s,t,-s-t$. Put $h(s,t)=s+t$ and absorb the sign of the third path into
$\sigma_e$: on the first two paths $\sigma_e=\varepsilon_e$, and on the
third $\sigma_e=-\varepsilon_e$. Thus
\begin{equation}\label{eq:theta-arc-form}
 x_e=v_e+\sigma_e\eta_e(s,t),\qquad
 \eta_e\in\{s,t,h\},\qquad h=s+t.
\end{equation}
The symbol $h$ here is a linear coordinate form, not the number of states.
Signed interval intersection as in \eqref{eq:path-intervals} gives
\begin{equation}\label{eq:theta-base}
 T_B=\{(s,t):L_s\le s\le U_s,\ L_t\le t\le U_t,
                       \ L_h\le s+t\le U_h\}.
\end{equation}
A cycle has the analogous scalar form $x_e=v_e+\sigma_es$ and base
interval $[L_s,U_s]$. Read the aggregate coordinates $s(x),t(x)$ from
one representative arc on each of the first two paths; then
$h(x)=s(x)+t(x)$. These representative arcs are distinct.

Fix a block and use its local states $I_B$ and weights from
\eqref{eq:merged-weight}. Suppress $B$ in the local formulas below, so
$\lambda_j=\lambda_{B,j}$, including the merged state $*$. Define
$\widetilde q_{ej}=\sigma_e(z_{ej}-v_ey_j)$ for observations in this
block. Initialize and tighten each coordinate interval by
\begin{equation}\label{eq:theta-tightened-intervals}
 \ell_{\eta j}=\max\bigl(\{\lambda_jL_\eta\}\cup
        \{\widetilde q_{ej}:(e,j)\in\Obs,\ \eta_e=\eta\}\bigr),\qquad
 u_{\eta j}=\min\bigl(\{\lambda_jU_\eta\}\cup
        \{\widetilde q_{ej}:(e,j)\in\Obs,\ \eta_e=\eta\}\bigr).
\end{equation}
For $*$ there are no observations. Multiple observations on a path must
agree: they all enter both endpoint lists. At zero weight the initial
interval is the singleton zero, so feasibility forces every observed
coordinate and every state deviation to be zero.

The cycle state domain is $[\ell_{sj},u_{sj}]$. The theta state domain is
\begin{equation}\label{eq:theta-state}
 Q_j=\{(s,t):\ell_{sj}\le s\le u_{sj},\
       \ell_{tj}\le t\le u_{tj},\
       \ell_{hj}\le s+t\le u_{hj}\}.
\end{equation}
By \cref{thm:block-state-reduction}, block membership asks whether the
aggregate coordinate belongs to the Minkowski sum of these local state
domains. This reduction holds at every simplex boundary and for
lower-dimensional state domains.

\begin{lemma}[Theta feasibility and tight supports]\label{lem:theta-supports}
The state domain $Q_j$ is nonempty if and only if
\begin{equation}\label{eq:theta-nonempty}
 \ell_{\eta j}\le u_{\eta j}\quad(\eta=s,t,h),\qquad
 \ell_{sj}+\ell_{tj}\le u_{hj},\qquad
 \ell_{hj}\le u_{sj}+u_{tj}.
\end{equation}
When it is nonempty, its attained minimum and maximum of each coordinate
form are
\begin{equation}\label{eq:theta-six-supports}
\begin{array}{ll}
 a_{sj}=\max\{\ell_{sj},\ell_{hj}-u_{tj}\},&
 b_{sj}=\min\{u_{sj},u_{hj}-\ell_{tj}\},\\[2pt]
 a_{tj}=\max\{\ell_{tj},\ell_{hj}-u_{sj}\},&
 b_{tj}=\min\{u_{tj},u_{hj}-\ell_{sj}\},\\[2pt]
 a_{hj}=\max\{\ell_{hj},\ell_{sj}+\ell_{tj}\},&
 b_{hj}=\min\{u_{hj},u_{sj}+u_{tj}\}.
\end{array}
\end{equation}
\end{lemma}
\begin{proof}
The rectangle of individual coordinate intervals is nonempty exactly
when its two intervals are nonempty. Its attainable sums fill
$[\ell_{sj}+\ell_{tj},u_{sj}+u_{tj}]$. Requiring that interval to
intersect the nonempty sum interval gives precisely
\eqref{eq:theta-nonempty}. For a fixed $s$, a feasible $t$ exists if and
only if $s\in[\ell_{sj},u_{sj}]$ and
$s\in[\ell_{hj}-u_{tj},u_{hj}-\ell_{tj}]$. This proves both tight
$s$ bounds and their attainment. Interchanging $s,t$ proves the next
pair, and intersecting the two attainable sum intervals proves the last.
\end{proof}

\begin{lemma}[Six supports suffice for planar sums]\label{lem:theta-minkowski}
For any finite family of nonempty domains of the form
\eqref{eq:theta-state}, including segments and points,
\begin{equation}\label{eq:theta-minkowski}
 \sum_jQ_j=
 \left\{(s,t):\sum_j a_{\eta j}\le\eta(s,t)
                    \le\sum_j b_{\eta j}\quad(\eta=s,t,h)\right\}.
\end{equation}
The empty sum is the singleton zero, represented by six zero supports.
\end{lemma}
\begin{proof}
Let $\mathcal N=\{\pm(1,0),\pm(0,1),\pm(1,1)\}$. Every edge of a
two-dimensional $Q_j$ has a normal in $\mathcal N$, because $Q_j$ is
defined by inequalities with those normals. A one-dimensional $Q_j$
also has a direction perpendicular to a member of $\mathcal N$: at a
relative interior point of the segment, at least one defining inequality
is tight and its boundary contains the segment. Otherwise that point
would have a two-dimensional neighborhood in $Q_j$.

For a linear objective with normal $n$, its maximizing face over a sum
is the sum of its maximizing faces over the summands. Indeed, objective
values add, and equality in the sum of the separate maximum values holds
exactly when every summand attains its maximum. If the sum is
two-dimensional, each exposed edge is therefore a sum of points and
parallel segments, with at least one nontrivial segment. Its direction
is a direction already occurring in a summand. Thus every edge normal of
the sum lies in $\mathcal N$, and its support values are the sums of
those in \cref{lem:theta-supports}. These supports describe the polygon.

If the sum is a nontrivial segment, each nonpoint summand is a segment
parallel to it. Its direction is one of the three permitted directions.
The two normals perpendicular to that direction enforce its supporting
line, and at least one of the other coordinate forms varies along that
line. The lower and upper supports of that form give exactly the segment
endpoints. If the sum is a point, the supports of $s$ and $t$ fix it.
This proves the same six-inequality description in all dimensions.
\end{proof}

The planar restriction matters. In higher dimension, an exposed facet
can be a sum of lower-dimensional faces with different directions, so the
preceding edge argument does not give a fixed list of facet normals.
The support and normal-fan identities used here are classical
\citep{GritzmannSturmfels1993}; the explicit formulas exploit these
particular block domains.

\begin{theorem}[Cycle/theta original-coordinate hull]\label{thm:theta-hull}
Suppose every cyclic block is a cycle or a theta block. The hull $\Hull$
is described by the original flow and simplex constraints, the bridge
equations, the cycle conditions
\begin{equation}\label{eq:cycle-hull}
 \ell_{sj}\le u_{sj}\quad(j\in I_B),\qquad
 \sum_{j\in I_B}\ell_{sj}\le s(x)\le\sum_{j\in I_B}u_{sj},
\end{equation}
and, for every theta block, the local conditions
\eqref{eq:theta-nonempty} and six aggregate conditions
\begin{equation}\label{eq:theta-aggregate}
 \sum_{j\in I_B}a_{\eta j}\le\eta(s(x),t(x))
                   \le\sum_{j\in I_B}b_{\eta j}\qquad(\eta=s,t,h).
\end{equation}
All endpoints are evaluated by \eqref{eq:theta-tightened-intervals}.
These expressions define a finite family of linear inequalities with
coefficients in $\{0,\pm1\}$ on $x,z$.
\end{theorem}
\begin{proof}
For intervals, nonemptiness is checked separately, and their Minkowski
sum has the summed lower and upper endpoints. The theta case follows
from \cref{lem:theta-supports,lem:theta-minkowski}. The exact equivalence
with the global hull then follows from
\cref{thm:block-state-reduction}.

Every lower endpoint in \eqref{eq:theta-tightened-intervals} is a maximum
of affine functions, and every upper endpoint is a minimum. These
properties also hold for the lower and upper supports in
\eqref{eq:theta-six-supports}: for example,
$\ell_{hj}-u_{tj}$ is a maximum of affine differences. Consequently
each displayed condition can be written as a maximum of finitely many
affine functions bounded above by zero. Requiring every affine branch
to be at most zero is an equivalent finite linear description.

In a support branch, a state's product contribution uses either one
coordinate type or two distinct types. Hence each observed product occurs
at most once. Different states have different product indices. A local
same-type interval comparison either uses two distinct observations or
cancels an identical one; other local comparisons use distinct types.
The aggregate $s(x),t(x)$ use separate representative arcs, and $h(x)$
uses their sum. Thus every flow or product coefficient is $0,1$, or
$-1$. Simplex coefficients contain the network data. The assertion
concerns this rational row scaling, with the normalization qualifications
given after \eqref{eq:compression-nonzeros}.
\end{proof}

\subsection{A linear-arithmetic oracle and compact recovery}\label{subsec:theta-oracle}
If any condition in \cref{thm:theta-hull} fails, select a maximizing
affine branch of its convex piecewise-affine left side after moving all
terms to one side. That branch is bounded above by zero at every hull
point and has the same positive violation at the candidate. Thus endpoint
and support comparisons produce a globally valid linear cut directly;
they do not merely produce a local linear approximation. Ties may be
resolved arbitrarily among active branches.

There are $a_B+1$ local states in block $B$, and
$\sum_Ba_B\le|\Obs|$. A constant number of endpoint and support
operations suffices for each cycle/theta state after its observations
have been scanned. Checking the original domains and all block conditions
therefore uses
\begin{equation}\label{eq:theta-oracle-complexity}
 O(|V|+|E|+m+|\Obs|)
\end{equation}
rational arithmetic operations, given observation indices grouped by
block and label. Graph blocks, paths, and reference flows can be
preprocessed in linear time. Sorting an ungrouped observation list adds
$O(|\Obs|\log|\Obs|)$ comparisons if required. These are arithmetic and
comparison counts, not constant-bit-cost claims. All evaluated rational
expressions and emitted cut coefficients have polynomial encoding length.

\begin{proposition}[Constructive cycle/theta recovery]\label{prop:theta-recovery}
For a rational hull point in the scope of \cref{thm:theta-hull}, a compact
rational decomposition into at most $m+1$ graph points can be computed
within \eqref{eq:theta-oracle-complexity}, with the same grouping proviso.
The compact representation stores the global weights, reference flow,
one default cycle-coordinate vector per block, and exceptions for observed block/label
pairs. Expanding all flow vectors can require $\Theta((m+1)|E|)$ output
entries.
\end{proposition}
\begin{proof}
Order the local states in a theta block and precompute suffix sums of
the six tight supports. Suppose the remaining aggregate is $r=(r_s,r_t)$,
with $r_h=r_s+r_t$, and the next state is $i$. Its permissible coordinate
lies in $Q_i\cap(r-\sum_{j>i}Q_j)$. By
\cref{lem:theta-minkowski}, this is another theta polygon with intervals
\begin{equation}\label{eq:theta-recovery-intersection}
 \max\left\{\ell_{\eta i},r_\eta-\sum_{j>i}b_{\eta j}\right\}
 \le\eta\le
 \min\left\{u_{\eta i},r_\eta-\sum_{j>i}a_{\eta j}\right\},
 \quad\eta=s,t,h.
\end{equation}
It is nonempty because the remaining aggregate belongs to the remaining
sum. In any nonempty theta polygon with interval endpoints $\ell,u$,
choose
\begin{equation}\label{eq:theta-point-choice}
 s=\max\{\ell_s,\ell_h-u_t\},\qquad
 t=\max\{\ell_t,\ell_h-s\}.
\end{equation}
By \cref{lem:theta-supports}, the chosen $s$ lies between its tight
bounds, in particular $s\le u_s$ and $s\le u_h-\ell_t$. Both
arguments defining $t$ are at most $u_t$ because
$s\ge\ell_h-u_t$. The sum is at least $\ell_h$; if $t=\ell_t$ it
is at most $u_h$, and otherwise it equals $\ell_h$. Thus the chosen
point is feasible. Subtract it from $r$ and continue. Empty suffixes
use zero supports. The scalar interval construction is the same
intersection-and-subtraction procedure in one dimension.

Normalize each positive-weight explicit state coordinate by its weight.
Normalize a positive-weight merged coordinate by its total weight and
store it as the block's default for all constituent global states.
If that weight is zero, the default is unused and may be the aggregate
block coordinate, which is feasible because $x\in\Flow$. An unobserved
block also uses its aggregate coordinate as the default. The refinement
proof of \cref{thm:block-state-reduction} combines these coordinates into
feasible global flows with weights $\lambda_0,\ldots,\lambda_m$; zero
weights are omitted. The number of exceptions is at most $|\Obs|$,
and all operations and stored coordinates obey the claimed bound.
Enumerating a dense flow for every positive state is a separate output cost.
\end{proof}

\begin{example}[Separate state hulls miss a joint restriction]\label{ex:joint-theta}
Consider three parallel arcs of unit capacity carrying unit total flow.
Let their first two aggregate flows be $X_s=X_t=1/3$, and let
$y_1=y_2=1/3$. Observe both first-two-arc products for both labels, and
set all four to zero. For either label alone, the weighted state flow
$(0,0,1/3)$ and its complementary flow $(1/3,1/3,0)$ give a feasible
one-label hull decomposition. Thus the candidate belongs to the
intersection of the two separately convexified one-label hulls.

Jointly, the two observed states would consume $2/3$ of the third arc,
whose aggregate flow is only $1/3$. The violated joint inequality is
\begin{equation}\label{eq:joint-theta-cut}
 z_{s1}+z_{t1}+z_{s2}+z_{t2}
       \ge X_s+X_t-(1-y_1-y_2).
\end{equation}
Its violation is $1/3$. It follows from the residual state's upper
bound on the sum of the first two flows, and is included in
\eqref{eq:theta-aggregate}. The example illustrates the known joint
disaggregation mechanism rather than a separate convexification principle.
\end{example}

\subsection{Any number of parallel paths}\label{subsec:parallel-paths}
Now suppose each cyclic block, apart from independent loops, consists of
$k\ge2$ internally vertex-disjoint paths between two terminals. The
number $k$ is unrestricted. Orient all path traversals from the first
terminal to the second, and use the signs $\varepsilon_e$ of
\eqref{eq:path-deviation}. The core domain is
\begin{equation}\label{eq:parallel-base}
 L_i\le\delta_i\le U_i\quad(i\in[k]),\qquad
                 \sum_{i=1}^k\delta_i=0.
\end{equation}
Again write $\lambda_j=\lambda_{B,j}$ for local states, including $*$.
For each local state $j\in I_B$, intersect its scaled path bounds with
all observations on that path:
\begin{equation}\label{eq:parallel-endpoints}
 \ell_{ij}=\max\bigl(\{\lambda_jL_i\}\cup
           \{q_{ej}:(e,j)\in\Obs,\ e\in i\}\bigr),\qquad
 u_{ij}=\min\bigl(\{\lambda_jU_i\}\cup
           \{q_{ej}:(e,j)\in\Obs,\ e\in i\}\bigr).
\end{equation}
Here $q_{ej}$ uses \eqref{eq:aggregate-path-expression}, without the
theta-specific sign change. Let $\bar\delta_i=s_i(x)$ be the aggregate
path deviations. Exact block membership is equivalent to finding a
matrix $w=(w_{ij})$ such that
\begin{equation}\label{eq:parallel-matrix}
 \ell_{ij}\le w_{ij}\le u_{ij},\qquad
 \sum_iw_{ij}=0\quad(j\in I_B),\qquad
 \sum_jw_{ij}=\bar\delta_i\quad(i\in[k]).
\end{equation}
This follows directly from \cref{thm:block-state-reduction}.

\begin{theorem}[Complete parallel-path subset description]\label{thm:parallel-hull}
In this graph class, $\Hull$ is described by the original domains, bridge
equations, scalar loop conditions, and the following conditions for every
parallel-path block:
\begin{equation}\label{eq:parallel-local}
 \ell_{ij}\le u_{ij}\quad(i,j),\qquad
 \sum_i\ell_{ij}\le0\le\sum_i u_{ij}\quad(j\in I_B),
\end{equation}
and, for every subset $S\subseteq[k]$,
\begin{equation}\label{eq:parallel-subset}
 \sum_{i\in S}\bar\delta_i\le
 \sum_{j\in I_B}\min\left\{\sum_{i\in S}u_{ij},
                           -\sum_{i\notin S}\ell_{ij}\right\}.
\end{equation}
After expanding affine branches, this is a finite linear description
whose coefficients on $x,z$ lie in $\{0,\pm1\}$, with the same
rational-scaling qualification as \cref{thm:theta-hull}.
\end{theorem}
\begin{proof}
The local conditions are necessary for \eqref{eq:parallel-matrix}.
In a zero-sum column, the sum on $S$ is bounded by its upper endpoints
on $S$ and by the negative lower endpoints on its complement. Summing
over columns proves \eqref{eq:parallel-subset}.

For sufficiency, shift all entries by their lower endpoints and define
\begin{equation}\label{eq:transport-shift}
 g_{ij}=w_{ij}-\ell_{ij},\quad C_{ij}=u_{ij}-\ell_{ij},\quad
 r_i=\bar\delta_i-\sum_j\ell_{ij},\quad c_j=-\sum_i\ell_{ij}.
\end{equation}
We seek a nonnegative matrix $g\le C$ with row sums $r$ and column
sums $c$. The local conditions give $C_{ij}\ge0$ and $c_j\ge0$.
For each $i$, apply \eqref{eq:parallel-subset} to $[k]\setminus\{i\}$.
Since $\sum_i\bar\delta_i=0$, this gives
\[
 -\bar\delta_i\le\sum_j
 \min\left\{\sum_{h\ne i}u_{hj},-\ell_{ij}\right\}
 \le-\sum_j\ell_{ij},
\]
so $r_i\ge0$ as well. Moreover $R:=\sum_ir_i=\sum_jc_j$.

Build a directed network with a source, one node for each row, one node
for each column, and a sink. Source-to-row capacities are $r_i$,
row-to-column capacities are $C_{ij}$, and column-to-sink capacities are
$c_j$. A flow of value $R$ saturates all source and sink arcs and yields
exactly the desired matrix. Fixing the source-side row set $S$, the
minimum cut over the side choices of column nodes has capacity
\begin{equation}\label{eq:transport-cut}
 \sum_{i\notin S}r_i+
       \sum_j\min\left\{\sum_{i\in S}C_{ij},c_j\right\}.
\end{equation}
Indeed, a sink-side column pays its incoming capacities from $S$, while
a source-side column pays $c_j$; the choices are independent. This cut
has capacity at least $R$ if and only if
\[
 \sum_{i\in S}r_i\le\sum_j
            \min\left\{\sum_{i\in S}C_{ij},c_j\right\}.
\]
Adding $\sum_{i\in S,j}\ell_{ij}$ to both sides converts this
inequality exactly into \eqref{eq:parallel-subset}. Hence every cut has
capacity at least $R$. The maximum-flow/minimum-cut theorem
\citep{FordFulkerson1956Flow} gives the desired flow, including the case
$R=0$, when all targets vanish. Undo the shift to obtain
\eqref{eq:parallel-matrix}, then refine local states and combine blocks
by \cref{thm:block-state-reduction}.

The right side of \eqref{eq:parallel-subset} is a minimum of finitely
many affine expressions: $u_{ij}$ and $-\ell_{ij}$ have this property,
and it is preserved by sums and minima. Branch expansion therefore
gives an equivalent finite linear family. In a branch for one state,
either one upper endpoint from each path in $S$ is used, or one negated
lower endpoint from each complementary path is used. Each observation
belongs to only one path and one state, so its coefficient has magnitude
at most one. The local conditions have the same property, with possible
cancellation of a repeated observation in a same-interval comparison.
The aggregate left side uses distinct representative arcs. This proves
the coefficient assertion.
\end{proof}

For $k=3$, the six nonempty proper subsets are the supports in
\eqref{eq:theta-six-supports}, after setting
$(\delta_1,\delta_2,\delta_3)=(s,t,-h)$. The empty and full subsets
are redundant under \eqref{eq:parallel-local}. For larger $k$, the
subset family may be exponential. Its completeness comes from the
transportation projection, not from an extension of the planar edge
argument. This is an application of classical transportation feasibility
\citep{FordFulkerson1956Transportation}; the network-Cayley projection
of \citet[Sections~5.7 and~5.9]{KisHorvath2022} is a direct predecessor.
The lower-bound shift and transportation cut projection are explicit in their
Section~5.9 (Proposition~22 and equations~(30)--(31)).

\subsection{Separation by flow and reconstruction}\label{subsec:parallel-oracle}
The theorem supplies an oracle without enumerating the subsets. Check
the original domains, bridge equations, and local conditions first.
If a shifted row target $r_i$ in \eqref{eq:transport-shift} is negative,
the valid condition
\begin{equation}\label{eq:negative-row-cut}
 \bar\delta_i\ge\sum_j\ell_{ij}
\end{equation}
separates the point; its validity was proved from the complementary
subset above. Otherwise all transportation capacities and targets are
nonnegative. Compute a maximum flow. Value $R$ proves block membership
and returns its state matrix. A smaller value gives a minimum cut;
its source-side row set violates \eqref{eq:parallel-subset}, because
minimizing column sides can only decrease the cut capacity.

To return an original-coordinate linear inequality, select the active
minimum alternatives and active affine endpoint expressions in the
violated subset condition. Their sum is an affine majorant of its
concave right side and agrees with it at the candidate. Hence the
resulting affine upper bound is globally valid and strictly violated.
Local violations and \eqref{eq:negative-row-cut} are handled by their
active affine branches in the same way.

For a block with $k_B$ paths and $a_B$ observed labels, the transportation
network has
\begin{equation}\label{eq:transport-size}
 k_B+a_B+3\text{ nodes},\qquad
 k_B(a_B+1)+k_B+a_B+1\text{ arcs}.
\end{equation}
Graph and endpoint preprocessing, target construction, and cut extraction
use
\begin{equation}\label{eq:parallel-build-complexity}
 O\left(|V|+|E|+m+|\Obs|+\sum_B k_B(a_B+1)\right)
\end{equation}
arithmetic operations after observation grouping, followed by one maximum
flow computation per parallel-path block. A polynomial-time algorithm
such as shortest augmenting paths \citep{EdmondsKarp1972} gives polynomial
bit complexity for rational data. One way to see the encoding bound is
to clear the finitely many capacity denominators with a common multiple;
its encoding length is at most their summed encoding lengths. Flow
updates then use integers bounded by the total capacity. An arbitrary
augmenting-path rule is not needed for this complexity claim. When
$k_B$ grows, \eqref{eq:parallel-build-complexity} and the flow solves
replace the linear cycle/theta bound.

The returned matrices provide exact rational block state deviations.
Normalize positive-weight columns and store the normalized merged column
as the default for its constituent global states. Zero-weight columns
vanish by their original scaled path bounds, and no division by zero
occurs. Together with the reference flow and block paths, the resulting
compact decomposition has storage of order
\eqref{eq:parallel-build-complexity}; expanding every dense state-flow
vector can still require $\Theta((m+1)|E|)$ entries.

The cycle rank of a parallel-path block is $k_B-1$, so this theorem
permits unbounded rank within a block. It does not cover every
series--parallel block. In particular, several serial parallel gadgets
inside one outer bypass can have interacting internal state profiles
that are not represented by the single transportation matrix
\eqref{eq:parallel-matrix}. The present results supply explicit
structural oracles and unit-coefficient descriptions; generic polynomial
separation already follows from \cref{prop:disaggregation}.
